% Plantilla del informe de Data in the Wild (estructura basada en los informes de ejemplo 2021, 2022 y 2025).
% Compilar con: latexmk -pdf main.tex
\documentclass[11pt,a4paper,twocolumn]{article}

\usepackage{iftex}
\ifPDFTeX
    \usepackage[utf8]{inputenc}
    \usepackage[T1]{fontenc}
    \usepackage{newtxtext,newtxmath}
\else
    % XeLaTeX / LuaLaTeX
    \usepackage{newtxmath}
    \usepackage[no-math]{fontspec}
    \setmainfont{texgyretermes}[Extension=.otf, UprightFont=*-regular, BoldFont=*-bold, ItalicFont=*-italic, BoldItalicFont=*-bolditalic]
\fi
\usepackage[english]{babel}
\usepackage{microtype}
\usepackage[a4paper,margin=2.5cm,columnsep=0.6cm]{geometry}
\usepackage{titlesec}
\titleformat*{\section}{\large\bfseries}
\titleformat*{\subsection}{\normalsize\bfseries}
\usepackage{graphicx}
\usepackage{booktabs}
\usepackage{caption}
\usepackage{subcaption}
\usepackage{enumitem}
\usepackage{amsmath}
\usepackage{siunitx}
\usepackage{xcolor}
\usepackage{csquotes}
\usepackage[style=numeric,sorting=none,backend=biber]{biblatex}
\usepackage[colorlinks=true,linkcolor=blue!60!black,citecolor=green!40!black,urlcolor=blue!60!black]{hyperref}
\usepackage[nameinlink]{cleveref}

\addbibresource{references.bib}
\graphicspath{{figures/}}

% Marcas para el borrador: quitarlas antes de la entrega final
\newcommand{\todo}[1]{\textcolor{red}{[TODO: #1]}}

% Datos del título
\newcommand{\reporttitle}{Comparing Supermarket Prices Across Europe: A Carrefour Shopping Basket}
\newcommand{\coursename}{Data in the Wild: Wrangling and Visualising Data (Autumn 2026)}
\newcommand{\coursecode}{KSDWWVD1KU}
\newcommand{\githuburl}{https://github.com/antoniolandin/carrefour}

% Resumen centrado y en letra pequeña, como en los informes a dos columnas de 2021 y 2022
\renewenvironment{abstract}
    {\begin{center}\bfseries\abstractname\end{center}\vspace{-0.5\baselineskip}\begin{quote}\small}
    {\end{quote}}

\title{\textbf{\reporttitle}}
\author{
    \textbf{Antonio Cabrera Landín} \quad \textbf{Javier \todo{apellidos}} \quad \textbf{\todo{resto del grupo}} \\
    IT University of Copenhagen \\
    \texttt{\{\todo{email},\todo{email}\}@itu.dk} \\
    \url{\githuburl}
}
\date{\small \coursename, course code \coursecode}

\begin{document}

% ---------------------------------------------------------------------------
% Título a todo el ancho; el enunciado pide el enlace a GitHub en la primera página
% ---------------------------------------------------------------------------
\maketitle

% ---------------------------------------------------------------------------
% Resumen (2021 y 2022 lo tienen; 150-250 palabras)
% ---------------------------------------------------------------------------
\begin{abstract}
    \todo{Qué hemos hecho (scraping de Carrefour en varios países), con qué datos, cómo emparejamos productos y la conclusión principal en una o dos frases.}
\end{abstract}

% ---------------------------------------------------------------------------
\section{Note on AI usage}
% Sección obligatoria en el informe de 2025: herramientas usadas y para qué.
\todo{Herramientas de IA usadas y para qué (código, depuración, redacción, gráficos).}

\begin{itemize}
    \item \todo{Uso 1}
    \item \todo{Uso 2}
\end{itemize}

% ---------------------------------------------------------------------------
\section{Introduction}
% Contexto, por qué no existe un dataset así y pregunta(s) de investigación explícitas.
\todo{Contexto: comparación de precios de la cesta de la compra entre países de la UE y por qué Carrefour es un buen caso.}

Our project aims to answer the following research questions:
\begin{itemize}
    \item \textbf{RQ1:} \todo{Por ejemplo: ¿cuánto cuesta la misma cesta de la compra en Carrefour en distintos países?}
    \item \textbf{RQ2:} \todo{Por ejemplo: ¿se venden los mismos productos (mismo EAN) y a qué precio?}
\end{itemize}

All code and data are available on GitHub (see \cref{app:github}).

% ---------------------------------------------------------------------------
\section{Data collection}
% Los tres ejemplos detallan fuentes, método (scraping/API), problemas encontrados y tamaño final.

\subsection{Shopping basket}
\todo{Cómo se eligieron los 22 productos de cesta.csv y su traducción a cada idioma.}

\subsection{Carrefour Spain}
\todo{Páginas de catálogo, bloque de analítica \texttt{impressions} con el EAN, tope de 1008 productos por listado.}

\subsection{Carrefour Italy}
\todo{Búsqueda con Playwright y JSON de cada tarjeta de producto.}

\subsection{Challenges}
\todo{Contenido generado con JavaScript, shadow DOM, bloqueos de Cloudflare, productos agotados.}

\begin{table}[ht]
    \centering
    \small
    \begin{tabular}{lrrr}
        \toprule
        Country & Products & EANs & Basket \\
        \midrule
        Spain   & \todo{}  & \todo{}     & \todo{} / 22 \\
        Italy   & \todo{}  & \todo{}     & \todo{} / 22 \\
        \bottomrule
    \end{tabular}
    \caption{\todo{Resumen de los datos recogidos por país: productos, EAN únicos y líneas de la cesta encontradas.}}
    \label{tab:datasets}
\end{table}

% ---------------------------------------------------------------------------
\section{Data processing}

\subsection{Schema harmonization}
\todo{Formato común para todos los países: campos, tipos, moneda.}

\subsection{Data quality}
\todo{Validación de EAN (dígito de control), códigos internos de frescos, duplicados, campos vacíos, nombres sin decimales.}

\subsection{Product matching across countries}
\todo{Emparejamiento por EAN, por nombre (marca + formato) y productos equivalentes de marca propia.}

% ---------------------------------------------------------------------------
\section{Annotation}
% Los tres ejemplos anotan datos a mano y miden el acuerdo entre anotadores.
\todo{Qué se anota (p. ej. qué producto corresponde a cada línea de la cesta), guía de anotación y acuerdo entre anotadores \cite{cohen1960}.}

% ---------------------------------------------------------------------------
\section{Results and discussion}
% En 2025: análisis exploratorio con muchas figuras; en 2021/2022: resultados y discusión separados.

\subsection{Exploratory data analysis}
\todo{Distribución de precios, cobertura de la cesta por país.}

% figure ocupa una columna; figure* (o table*) ocupa las dos y va arriba de una página
\begin{figure}[ht]
    \centering
    % \includegraphics[width=\linewidth]{precios_por_pais.pdf}
    \framebox[\linewidth]{\rule{0pt}{4cm}\todo{figura}}
    \caption{\todo{Pie de figura.}}
    \label{fig:prices}
\end{figure}

\subsection{Basket price comparison}
\todo{Respuesta a RQ1 (ver \cref{fig:prices}).}

\subsection{Shared products}
\todo{Respuesta a RQ2.}

% ---------------------------------------------------------------------------
\section{Limitations and ethical considerations}
% 2022 tiene una sección de limitaciones; aquí también encaja la ética del scraping.
\todo{Alcance distinto por país (catálogo frente a búsqueda), precios de un único día y código postal, términos de uso y bloqueos de la web.}

% ---------------------------------------------------------------------------
\section{Conclusion}
\todo{Respuesta breve a cada pregunta de investigación y trabajo futuro.}

\printbibliography[heading=bibintoc]

% ---------------------------------------------------------------------------
\appendix
\section{Appendix}

\subsection{GitHub}
\label{app:github}
The code and the datasets are available at \url{\githuburl}.

\subsection{Annotation guidelines}
\todo{Guía de anotación con ejemplos y casos dudosos.}

\subsection{Datasheet for the dataset}
% Estructura de Datasheets for Datasets, usada en el informe de 2025.
Following \textcite{gebru2021datasheets}:
\begin{description}[style=nextline]
    \item[Motivation] \todo{}
    \item[Composition] \todo{}
    \item[Collection process] \todo{}
    \item[Preprocessing, cleaning and labelling] \todo{}
    \item[Uses] \todo{}
    \item[Distribution] \todo{}
    \item[Maintenance] \todo{}
\end{description}

\end{document}
